\chapter{Paracentesis}

\section*{Overview}
Paracentesis is a procedure that removes fluid from the abdominal cavity through a needle or catheter, performed for diagnosis and treatment of ascites.

\section*{Alternative Names}
Abdominal paracentesis; abdominal tap; ascites drainage; therapeutic paracentesis

\section*{Principle}
Fluid accumulates in the peritoneal cavity in cirrhosis, malignancy, infection, and heart failure. A needle or catheter is inserted through the abdominal wall to sample or drain this fluid for analysis or symptom relief.

\section*{History}
Paracentesis has been performed since antiquity, and its role in diagnosing spontaneous bacterial peritonitis and relieving massive ascites was established in the 20th century.

\section*{Why It Is Done}
Performed for new or worsening ascites to determine its cause, to evaluate suspected spontaneous bacterial peritonitis, and to relieve respiratory or abdominal discomfort from large-volume ascites.

\section*{Is This a Primary or Secondary Test or Procedure}
Primary diagnostic and therapeutic procedure for ascites.

\section*{How It Is Performed}
Performed at the bedside. The patient lies supine with the head elevated, the puncture site is identified (usually the left or right lower quadrant, avoiding surgical scars and the epigastric vessels), the skin is prepped and anesthetized, and a needle or catheter is advanced through the abdominal wall. Fluid is aspirated or drained into collection bags.

\section*{Preparation}
Coagulation assessment and platelet count, bladder emptying, and ultrasound may be used to localize fluid in difficult cases.

\section*{Contraindications and Precautions}
Relative contraindications include severe coagulopathy, bowel obstruction, and overlying skin infection. Precaution: avoid the inferior epigastric artery and surgical scars where bowel may be adherent.

\section*{Risks}
Bleeding, infection, bowel perforation, and post-procedure hypotension with large-volume drainage; albumin is often given to prevent circulatory dysfunction.

\section*{Aftercare and Recovery}
The patient is observed briefly and monitored for bleeding or hypotension. Fluid is sent for cell count, culture, albumin, and cytology as indicated.

\section*{Outcome and Follow-up}
The serum-ascites albumin gradient (SAAG) distinguishes portal hypertension from other causes; an elevated neutrophil count indicates spontaneous bacterial peritonitis.

\section*{Expected Results}
Clear fluid with a low cell count and no organisms, and relief of the symptoms of ascites.

\section*{Follow-up}
Repeat paracentesis for recurrent symptoms, serial evaluation of the underlying liver or malignant disease, and monitoring for infection.

\section*{When to Seek Medical Care}
Follow the procedure team's specific instructions. Seek urgent medical assessment for severe or worsening pain, uncontrolled bleeding, fever or signs of infection, trouble breathing, chest pain, fainting, new weakness or confusion, inability to urinate, or any rapidly worsening symptom. The appropriate warning signs depend on the procedure and the patient's medical condition.

\section*{References}
\begin{enumerate}
  \item MedlinePlus. Paracentesis. U.S. National Library of Medicine; 2025.
  \item Current Medical Diagnosis and Treatment. McGraw-Hill; 2025.
\end{enumerate}

\clearpage
